% !PRINTED 783
\weberpage{783}
%<<CONT
(worin der imaginäre Theil des Logarithmus zwischen $\pm \pi i$ zu nehmen ist), so folgt, wenn der imaginäre Theil von $\log Q(s)$ passend bestimmt wird,

\weq{21}{\log Q(s) = \Sigma \frac{\chi(p)}{p^s} + \frac{1}{2}\Sigma \frac{\chi(p^2)}{p^{2s}} + \frac{1}{3}\Sigma \frac{\chi(p^3)}{p^{3s}} + \cdots}

Wenn nun $A$ irgend eine der Classen (2) und $a$ eine Zahl aus $A$ bedeutet, so ist nach \S.\,11, (6) dieses Bandes

\weq{22}{\overset{\chi}{\Sigma} \chi(A) \chi(n) = \Sigma \chi(an) = \mu \text{ oder } = 0,}

je nachdem $an \equiv 1$ oder nicht $\equiv 1$ nach dem Modul $m$ ist, d.\,h. je nachdem $n$ in der Classe $A^{-1}$ enthalten ist oder nicht enthalten ist. Wenn wir also die Formel (21) mit $\chi(A)$ multipliciren und in Bezug auf $\chi$ summiren, so folgt

\weq{23}{\frac{1}{\mu} \overset{\chi}{\Sigma} \chi(A) \log Q(s) = \Sigma \frac{1}{p^s} + \frac{1}{2}\Sigma \frac{1}{p^{2s}} + \frac{1}{3}\Sigma \frac{1}{p^{3s}} + \cdots,}

wo sich rechts die erste, zweite, dritte, \ldots\ Summe auf alle Primzahlen $p$ bezieht, deren erste, zweite, dritte, \ldots\ Potenz in $A^{-1}$ enthalten ist.

Der zweite Theil dieser Summe
\[
R(s) = \frac{1}{2}\Sigma \frac{1}{p^{2s}} + \frac{1}{3}\Sigma \frac{1}{p^{3s}} + \cdots
\]
wird vergrössert, wenn man jede seiner Theilsummen auf \gesperrt{alle} Primzahlen $p$ erstreckt, und demnach ist
\[
R(s) < \frac{1}{2}\Sigma \left(\frac{1}{p^{2s}} + \frac{1}{p^{3s}} + \cdots\right) = \frac{1}{2}\Sigma \frac{1}{p^s (p^s - 1)},
\]
also um so mehr:
\[
R(s) < \frac{1}{2} \sum_{2,\infty}^{n} \frac{1}{n^{2s} (1 - n^{-s})},
\]
oder, da $1 - n^{-s} > \frac{1}{2}$ ist:
\[
R(s) < \sum_{2,\infty}^{n} \frac{1}{n^{2s}},
\]
und hieraus schliesst man, dass $R$ für $s = 1$ endlich bleibt (\S.\,191, 2.).

Wenn wir nun in der Formel (23),
\[
\frac{1}{\mu} \overset{\chi}{\Sigma} \chi(A) \log Q(s) = \Sigma \frac{1}{p^s} + R(s),
\]
%CONT>>
